% GENERATED by scripts/build-lecture-slides.py; do not edit. Portable profile: XeTeX, theme and fonts from the core.
\makeatletter\IfFileExists{core/beamer/beamerthememx.sty}{\def\mxcorepath{core/}\def\input@path{{core/beamer/}{core/}}}{\def\mxcorepath{../core/}\def\input@path{{../core/beamer/}{../core/}}}\makeatother
\documentclass[aspectratio=169,10pt,t]{beamer}
\usetheme{mx}
\usepackage{graphicx,booktabs}
\usepackage[french]{babel}
\setlength{\parskip}{.45em}
\hypersetup{pdfauthor={},pdfsubject={Lecture slides},pdfkeywords={},pdftitle={Statistiques}}
\providecommand{\N}{\mathbb{N}}
\providecommand{\Z}{\mathbb{Z}}
\providecommand{\D}{\mathbb{D}}
\providecommand{\Q}{\mathbb{Q}}
\providecommand{\R}{\mathbb{R}}
\providecommand{\C}{\mathbb{C}}
\providecommand{\vect}[1]{\overrightarrow{#1}}
\providecommand{\vi}{\vec{\imath}}
\providecommand{\vj}{\vec{\jmath}}
\providecommand{\vk}{\vec{k}}
\providecommand{\norm}[1]{\left\lVert #1\right\rVert}
\providecommand{\vc}[2]{\begin{pmatrix}#1\\#2\end{pmatrix}}
\providecommand{\deter}[2]{\det\left(#1,#2\right)}
\providecommand{\abs}[1]{\left\lvert #1\right\rvert}
\providecommand{\dd}{\,\mathrm{d}}
\providecommand{\e}{\mathrm{e}}
\providecommand{\Card}{\operatorname{Card}}
\providecommand{\E}{\operatorname{E}}
\providecommand{\V}{\operatorname{V}}
\providecommand{\degre}{^{\circ}}
\providecommand{\un}[1]{\,\mathrm{#1}}
\providecommand{\sci}[2]{#1\times10^{#2}}
\providecommand{\pgcd}{\operatorname{PGCD}}
\providecommand{\ppcm}{\operatorname{PPCM}}
\title{Statistiques}
\subtitle{Cours et exemples}
\date{}
\author{}
\begin{document}
\mxtitleframe
\begin{frame}{Séries statistiques à une variable}
Considérons le tableau suivant qui donne les âges des étudiants d'une classe
                  de terminale comptant 43 élèves : 

\begin{center}\begin{tabular}{lllllll}
Ages & 15 & 16 & 17 & 18 & 19 & 20 \\ 
Nombre d'élèves & 2 & 6 & 17 & 14 & 3 & 1 \\ 
\end{tabular}\end{center}
\end{frame}
\begin{frame}{Séries statistiques à une variable (2)}
\begin{itemize}
\item L'ensemble des 43 élèves constitue la \emph{population statistique}.

\item Chaque élève est un \emph{individu} ou \emph{unité statistique}.

\item Le \emph{caractère statistique} étudié dans ce tableau est l'âge.

\end{itemize}
\end{frame}
\begin{frame}{Séries statistiques à une variable (3)}
\begin{itemize}
\item Il s'agit d'un \emph{caractère quantitatif}, c'est-à-dire qu'il peut s'exprimer en nombre.
                              Un caractère qui ne peut pas s'exprimer en nombre est dit \emph{qualitatif}.

\item Le caractère étudié ici prend un nombre fini de valeurs, on dit qu'il est \emph{discret}.

                              Un caractère qui peut prendre n'importe quelle valeur réel dans un intervalle 

                              est dit \emph{continu}.

\end{itemize}
\end{frame}
\begin{frame}{Séries statistiques à une variable (4)}
\begin{itemize}
\item Un caractère quantitatif est aussi appelé \emph{variable}.

\item Pour étudier une variable continue, on partage l'intervalle dans lequel elle prend valeurs en 

                              un nombre fini d'intervalle qu'on appelle \emph{classe}.
                        
\item Les \emph{modalités} d'une variable discrète sont les valeurs que peut prendre la variable.

\end{itemize}
\end{frame}
\begin{frame}{Séries statistiques à une variable (5)}
\begin{itemize}
\item L'\emph{effectif} d'une modalité est le nombre d'individus présentant cette modalité.

\item La \emph{fréquence} d'une modalité est le 

                              rapport entre l'effectif de cette modalité et l'effectif total.

\end{itemize}
\end{frame}
\begin{frame}{Séries statistiques à une variable (6)}
\begin{itemize}
\item Dans le cadre d'une variable quantitative, l'\emph{effectif cumulé}

                              d'une modalité \(a\) est la somme des effectifs des modalités inférieures ou égales à \(a\).
                              De même, la \emph{fréquence cumulée} d'une modalité \(a\) est la somme des fréquences

                              des modalités inférieures ou égales à \(a\).

\end{itemize}
\end{frame}
\begin{frame}{Séries statistiques à une variable (7)}
example
                        On donne la série statistique du 
                        nombre de frères et sœurs des étudiants d'une classe de terminale :

\begin{center}\begin{tabular}{llllllllll}
1 & 2 & 0 & 1 & 1 & 0 & 2 & 3 & 2 & 3 \\ 
3 & 0 & 4 & 1 & 2 & 2 & 4 & 2 & 0 & 3 \\ 
1 & 2 & 0 & 2 & 5 & 1 & 2 & 2 & 3 & 0 \\ 
\end{tabular}\end{center}
\end{frame}
\begin{frame}{Séries statistiques à une variable (8)}
Ces données peuvent être représentées de manière plus organisée par le tableau suivant :
                        table[H]
\end{frame}
\begin{frame}{Séries statistiques à une variable (9)}
\begin{center}\begin{tabular}{lllll}
Valeur & Effectif & Fréquence & Effectifs & Fréquence \\ 
(\(x_i\)) & (\(n_i\)) & (\(f_i\)) & cumulés & cumulés \\ 
0 & 6 & 0,2 & 6 & 0,2 \\ 
1 & 6 & 0,2 & 12 & 0,4 \\ 
2 & 10 & 0,33 & 22 & 0,73 \\ 
3 & 5 & 0,17 & 27 & 0,90 \\ 
4 & 2 & 0,07 & 29 & 0,97 \\ 
5 & 1 & 0,03 & 30 & 1,0 \\ 
\end{tabular}\end{center}
\end{frame}
\begin{frame}{Séries statistiques à une variable (10)}
Tableau représentant le
                                 nombre de frères et sœurs des étudiants d'une classe de terminale.
                        
                        table
                  example
\end{frame}
\begin{frame}{Séries statistiques à une variable}
Il est souvent préférable de visualiser les données par une représentation graphique.
\end{frame}
\begin{frame}{Séries statistiques à une variable}
\textbf{Diagramme en bâtons}

 Les diagrammes en bâtons sont utilisés pour représenter des variables discrètes.
                        La série statistique précédente peut par exemple être 
                        représenté par un diagramme en bâtons (Figure fig:diagramme en batons).
\end{frame}
\begin{frame}{Séries statistiques à une variable}
\textbf{Histogramme}

 Considérons le tableau tab:variable continue.
                        On a regroupé les données par classe d'amplitude 5 cm. 
                        table[h!]
\end{frame}
\begin{frame}{Séries statistiques à une variable (2)}
\begin{center}\begin{tabular}{lll}
Classes & Effectifs & Centre \\ 
[150 ; 155[ & 2 & 152.5 \\ 
[155 ; 160[ & 4 & 157.5 \\ 
[160 ; 165[ & 5 & 162.5 \\ 
[165 ; 170[ & 8 & 167.5 \\ 
[170 ; 175[ & 4 & 172.5 \\ 
[175 ; 180[ & 3 & 177.5 \\ 
[180 ; 185[ & 2 & 182.5 \\ 
[185 ; 190[ & 1 & 187.5 \\ 
\end{tabular}\end{center}
\end{frame}
\begin{frame}{Séries statistiques à une variable (3)}
Tableau présentant les tailles en centimètres des élèves d'une classe de terminale.
                        
                        table

Cette série peut être représentée graphiquement par le diagramme fig:histogramme.
\end{frame}
\begin{frame}{Séries statistiques à une variable}
Dans un histogramme, 
                              il est souvent préférable de travailler avec des classes d'amplitudes constantes.
                              Dans le cas où les amplitudes des classes sont différentes,
                              l'effectif doit être proportionnel à l'aire (et non à la hauteur).
\end{frame}
\begin{frame}{Séries statistiques à une variable}
\textbf{Diagramme en secteur}

 Le tableau tab:religions indique les religions les plus pratiquées
                        ainsi que le nombre de pratiquants en 2010.

table[h!]
\end{frame}
\begin{frame}{Séries statistiques à une variable (2)}
\begin{center}\begin{tabular}{llll}
Religions & Nombre de & Pourcentages & Angle du secteur \\ 
 & pratiquants &  & correspondant \\ 
Christianisme & 2 420 & 36,24 & 130,5 \\ 
Islam & 1 800 & 26,95 & 97,04 \\ 
Sans religion & 1 150 & 17,22 & 61,99 \\ 
Hindouisme & 520 & 7,79 & 28 \\ 
Bouddhisme & 400 & 5,99 & 21,56 \\ 
Autres & 388 & 5,81 & 20,92 \\ 
\end{tabular}\end{center}
\end{frame}
\begin{frame}{Séries statistiques à une variable (3)}
Tableau indiquant les religions les plus pratiqués 
                                       ainsi que le nombre de pratiquants (en millions) en 2010.
                              
                        table
\end{frame}
\begin{frame}{Séries statistiques à une variable}
Ce tableau peut être représenté par un diagramme en secteur 
                        comme dans la figure fig:religions.
\end{frame}
\begin{frame}{Séries statistiques à une variable}
\begin{itemize}
\item Dans le cas d'une variable discrète,

                                          le \emph{mode} correspond à la valeur du caractère 
                                          ayant l'effectif le plus élevé.
                                    
\item Dans le cas d'une variable continue,

                                          la \emph{classe modale} désigne la classe 
                                          qui correspond à l'effectif le plus élevé.

\end{itemize}
\end{frame}
\begin{frame}{Séries statistiques à une variable (2)}
\begin{itemize}
\item La \emph{médiane} est la valeur du caractère qui partage

                                          la série en deux groupes de même effectif.
                                          example
                                                
\begin{itemize}

                                                      \item On considère l'âge de dix personnes
                                                            rangées par ordre croissant :
                                                             
                                                                  15 
                                                                  16 
                                                                  16 
                                                                  17 
                                                                  17 
                                                                  18 
                                                                  18 
                                                                  18 
                                                                  19 
                                                                  19 .

                                                            On peut prendre comme médiane le milieu de \(]17 ; 18[\)
                                                            soit \(17,5\) ; c'est-à-dire 17 ans et 6 mois.
                                                            Ce n'est pas nécessaire ; 
                                                            toute valeurs entre 17 et 18 divise la série en deux.
                                                      
\item Considérons le tableau tab:veufs, 
                                                            montrant le nombre de veufs
                                                            dans une ville classé par catégorie d'âge.
                                                            table[H]

\begin{center}\begin{tabular}{lll}
Classes (\(c_i\)) & Effectifs (\(n_i\)) & Effectifs cumulé croissant (\(n_i\uparrow\)) \\ 
- de 30 & 15 & 15 \\ 
[30 ; 40[ & 48 & 63 \\ 
[40 ; 50[ & 169 & 232 \\ 
[50 ; 60[ & 432 & 664 \\ 
[60 ; 70[ & 858 & 1522 \\ 
[70 ; 80[ & 1388 & 2910 \\ 
[80 ; 90[ & 808 & 3718 \\ 
90 et + & 118 & 3836 \\ 
\end{tabular}\end{center}

                                                            Tableau représentant le nombre de veufs dans une ville.
                                                            
                                                            table
                                                            
                                                            Le total des effectifs est 3 836 et la médiane qui correspond
                                                            à un effectif \(\dfrac{N}{2} = 1918\) est atteinte dans la classe 
                                                            [70 ; 80[.
                                                            Comme l'effectif cumulé correspondant à la classe moins de 70
                                                            est de 1522 et celui de la classe moins de 80 est de 2910,
                                                            on interpole linéairement dans la classe [70 ; 80[.
                                                            La médiane est alors
                                                            
                                                              [.27]

                                                                Équation de la droite oblique :
                                                                \[
                                                                  y=\dfrac{\beta-\alpha}{b-a}(x-a)+\alpha.
                                                                \]
                                                                On en déduit
                                                                \[
                                                                  x=\dfrac{b-a}{\beta-\alpha}(y-\alpha)+a.
                                                                \]
                                                                Pour trouver \(x_m\), on prend \(y=m\) :
                                                                \[
                                                                  x_m=\dfrac{b-a}{\beta-\alpha}(m-\alpha)+a.
                                                                \]
                                                              
                                                             :
                                                            \[
                                                                  70 + \dfrac{80 - 70}{1388} \cdot (1918 - 1522)
                                                                   = 72,853\text{ ans}            
                                                            \]
                                                            soit 72 ans 10 mois et 7 jours.

\end{itemize}

                                          example

\end{itemize}
\end{frame}
\begin{frame}{Séries statistiques à une variable (3)}
\begin{itemize}
\item La \emph{moyenne arithmétique simple} est 

                                          \[
                                          \overline{x} = \dfrac{x_1 + x_2 + \cdots + x_n}{N}.
                                          \]

\end{itemize}
\end{frame}
\begin{frame}{Séries statistiques à une variable (4)}
\begin{itemize}
\item La \emph{moyenne arithmétique pondérée} est 

                                          \[
                                          \overline{x} = \dfrac{k_1 \cdot x_1 + k_2 \cdot x_2 + \cdots + k_n \cdot x_n}{N}.
                                          \]
                                          example
                                                
\begin{itemize}

                                                      \item Les notes de mathématiques d'un élève durant le premier 
                                                            trimestre sont : 8, 12 et 13.
                                                            La moyenne (simple) est donc \(\dfrac{8 + 12 + 13}{3} = 11\).
                                                      
\item Les notes d'un élève au premier trimestre sont représentés
                                                            dans le tableau tab:notes.
                                                            table[H]

\begin{center}\begin{tabular}{lll}
Matières & Notes/20 & Coefficients \\ 
Malagasy & 14 & 3 \\ 
Français & 13 & 2 \\ 
Anglais & 14 & 1 \\ 
Philosophie & 12 & 2 \\ 
Histo-géo & 11 & 2 \\ 
Mathématiques & 16 & 4 \\ 
Physique & 13 & 4 \\ 
SVT & 15 & 4 \\ 
EPS & 08 & 1 \\ 
\end{tabular}\end{center}

                                                                  Notes d'un élève en série scientifique.
                                                                  
                                                            table
                                                            
                                                            La moyenne est donc : 
                                                            \(\displaystyle{
                                                                  \dfrac{  3 \cdot 14 + 2 \cdot 13 + 1 \cdot 14
                                                                        + \cdots
                                                                        + 4 \cdot 15 + 1 \cdot 08 }
                                                                       {3 + 2 + 1 + \cdots + 4 + 1}
                                                            }\)

qui est égale à \(13,57\).

\end{itemize}

                                          example

\end{itemize}
\end{frame}
\begin{frame}{Séries statistiques à une variable (5)}
\begin{itemize}
\item Les \emph{quartiles} \(Q1\), \(Q2\) et \(Q3\) sont les trois valeurs 

                                          du caractère qui partagent la série en quatre groupes de même effectif.
                                          example
                                                On donne le tableau tab:quartiles.
                                                table[H]
                                                \[
                                                      \begin{array}{|c|c|c|c|}
                                                      \hline
                                                      x_i      & n_i      & n_i \uparrow & n_i \downarrow \\
                                                      \hline
                                                      7        & 4        & 4            & 27             \\
                                                      8        & 1        & 5            & 23             \\
                                                      9        & 2        & 7            & 22             \\
                                                      10       & 3        & 10           & 20             \\
                                                      11       & 3        & 13           & 17             \\
                                                      12       & 1        & 14           & 14             \\
                                                      13       & 5        & 19           & 13             \\
                                                      14       & 3        & 22           &  8             \\
                                                      15       & 4        & 26           &  5             \\
                                                      16       & 1        & 27           &  1             \\
                                                      \hline
                                                      \end{array}
                                                \]
                                                Tableau représentant des données inconnues.
                                                
                                                table
                                                On a \(N = 27\) et donc \(\dfrac{N}{4} = 6,75\).
                                                Comme \(Q1\) doit être une valeur de la série, on prend la valeur de rang \(7\)
                                                (premier entier supérieur ou égal à \(6,75\)) ; de même
                                                \(\dfrac{N}{2} = 13,5\) donne le rang \(14\) et \(\dfrac{3N}{4} = 20,25\) le rang \(21\).
                                                En lisant les effectifs cumulés croissants, on a alors :
                                                
\begin{itemize}

                                                      \item \(Q1 = 9\),
                                                      
\item \(Q2 = 12\) (c'est la médiane),
                                                      
\item \(Q3 = 14\).

\end{itemize}

                                          example

\%[]
                                               Avec la convention adoptée ici, et contrairement à la médiane, les quartiles doivent
                                                     nécéssairement être des valeurs de la série.
                                                     Dans le cas où \(N/4\) n'est pas un nombre entier,
                                                     on prend le plus petit entier supérieur (par exemple 7 pour 6,75).

\end{itemize}
\end{frame}
\begin{frame}{Séries statistiques à une variable (6)}
\begin{itemize}
\item L'\emph{interquartile} est \(Q3 - Q1\).

\item Les \emph{déciles} \(D1\), \(D2\), \ldots{} , \(D9\) sont les valeurs
                                          partageant la série en dix groupes de même effectif.

\end{itemize}
\end{frame}
\begin{frame}{Séries statistiques à une variable}
La variance est la moyenne des carrés des écarts à la moyenne, elle est définie par :

\[
                                          \sigma^2 = \dfrac{1}{N} \cdot \sum n_i \cdot (x_i - \overline{x})^2 .
                                    \]

                        L'écart-type est la racine carrée de la variance.
\end{frame}
\begin{frame}{Séries statistiques à une variable}
On considère le tableau tab:variance.
                              table[H]
\end{frame}
\begin{frame}{Séries statistiques à une variable (2)}
\begin{center}\begin{tabular}{llllll}
\(x_i\) & \(n_i\) & \(c_i\) & \(c_i - \overline{x}\) & \((c_i - \overline{x})^2\) & \(n_i \cdot (c_i - \overline{x})^2\) \\ 
[0;4[ & 4 & 2 & -8,76 & 73,7376 & 306,9504 \\ 
[4;8[ & 20 & 6 & -4,76 & 22,6576 & 453,152 \\ 
[8;12[ & 40 & 10 & -0,76 & 0,5776 & 23,104 \\ 
[12;16[ & 25 & 14 & 3,24 & 10,4976 & 262,44 \\ 
[16;20[ & 11 & 18 & 7,24 & 52,4176 & 576,5936 \\ 
Total : & 100 &  &  &  & 1622,24 \\ 
\end{tabular}\end{center}
\end{frame}
\begin{frame}{Séries statistiques à une variable (3)}
Notes des élèves dans une classe de terminale.
                                    
                              table
                              On a \(\overline{x} = \dfrac{1}{N} \cdot \sum n_i \cdot c_i\) et 
                              \(\sigma^2 = \dfrac{1622,24}{100} = 16,2224\).
                              Ainsi, \(\sigma = \sqrt{16,2224} = 4,0277\).
\end{frame}
\begin{frame}{Séries statistiques à une variable}
Plus l'écart-type est faible, plus la moyenne est représentative de la série.
\end{frame}
\begin{frame}{Séries statistiques à deux variables}
Soient \(x\) et \(y\) deux caractères quantitatifs d'une même population.
                  A chaque individu de la population, on associe un couple \((x_i ; y_i)\) 
                  où \(x_i\) et \(y_i\) sont les valeurs prises par les variables \(x\) et \(y\) respectivement.
                  
                  Dans un repère \((O ; \overrightarrow{\imath} , \overrightarrow{\jmath})\), l'ensemble des points \(M_i (x_i ; y_i)\)
                  est appelé \emph{nuage de points} associé à la série statistique.
                  
                  Soit une série statistique à deux variables \(x\) et \(y\) de moyenne \(\overline{x}\) et \(\overline{y}\).
                  Le point \(G\) de coordonnées \((\overline{x} ; \overline{y})\) est appelé \emph{point moyen}
                  du nuage de points associé à la série statistique.
                  
                  On appelle \emph{covariance} de \(x\) et \(y\) le nombre \(\sigma_{xy}\) défini par
                  
                  \(\displaystyle
                        \begin{aligned}
                        \sigma_{xy}
                           &=\dfrac{1}{n}\sum_{i=1}^{n}(x_i-\overline{x})(y_i-\overline{y}) \\
                           &=\dfrac{1}{n}\sum_{i=1}^{n}x_i y_i-\overline{x}\,\overline{y}.
                        \end{aligned}\)
                  
                  example
                        Soit la série statistique double suivante :
	                        \(x_i\) représente les salaires en milliers d'ariary proposés par diverses entreprises pour un poste,
	                        \(y_i\) indique le nombre de candidats qui se sont présentés pour ces emplois.
\end{frame}
\begin{frame}{Séries statistiques à deux variables (2)}
\begin{center}\begin{tabular}{llllll}
\(x_i\) & 44 & 45 & 46 & 47 & 48 \\ 
\(y_i\) & 10 & 13 & 17 & 19 & 21 \\ 
\end{tabular}\end{center}
\end{frame}
\begin{frame}{Séries statistiques à deux variables (3)}
.
                        On a \(\overline{x} = \dfrac{44 + 45 + 46 + 47 + 48}{5} = 46\)
                        et \(\overline{y} = \dfrac{10 +  13 + 17 + 19 + 21}{5} = 16\) ; 
	                        ainsi, le point moyen est \(G (46 ; 16)\).
                        
                        \[\begin{array}{rcl}
                        \sigma_{xy} & = & \dfrac{1}{5}
                                      [
                                            (44-46)(10 - 16) + (45-46)(13-16) + (46-46)(17-16) \\
                                    &   & + (47-46)(19-16) + (48-46)(21-16) ] \\
                                    & = & 5,6.
                        \end{array}\]

                  example
\end{frame}
\begin{frame}{Séries statistiques à deux variables}
Certaines séries statistiques représentent des variables qui sont liées.
                  Il est parfois possible de trouver une droite permettant d'approximer le nuage de points.
\end{frame}
\begin{frame}{Séries statistiques à deux variables}
Droite de régression de \(y\) en \(x\)

 La droite de régression obtenue par la méthode des moindres carrés passe par le point moyen ;
                        ainsi, elle a une équation de la forme \(y - \overline{y} = m (x - \overline{x})\) et le coefficient
                        directeur \(m\) est égal à \(\dfrac{\sigma_{xy}}{V(x)}\).
\end{frame}
\begin{frame}{Séries statistiques à deux variables}
Droite de régression de \(x\) en \(y\)

 Elle est donnée par l'équation \(x - \overline{x} = m' (y - \overline{y})\) et le coefficient
                        directeur \(m'\) est égal à \(\dfrac{\sigma_{xy}}{V(y)}\).
\end{frame}
\begin{frame}{Séries statistiques à deux variables}
\textbf{Corrélation linéaire}

 Le coefficient de corrélation linéaire sert à mesurer la relation entre les deux variables
                        \(x\) et \(y\).
                        Ce coefficient est donné par \(\displaystyle{r = \dfrac{\sigma_{xy}}{\sigma_x \cdot \sigma_y}}\).
\end{frame}
\begin{frame}{Séries statistiques à deux variables}
Un coefficient de corrélation proche de \(1\) ou \(-1\) signifie que 
                        les deux variables sont << fortement corrélées >>. 
	                        Comme \(\sigma_x\), \(\sigma_y\), \(V(x)\) et \(V(y)\) sont strictement positifs (aucune des deux variables n'est constante), 
                        les signes de \(m\), \(m'\) et \(r\) dépendent de celui de \(\sigma_{xy}\).
                              \[\begin{aligned}
                                    m\cdot m' & = \dfrac{\sigma_{xy}}{V(x)} \cdot \dfrac{\sigma_{xy}}{V(y)}        
                                                = \dfrac{(\sigma_{xy})^2}{\sigma_x^2 \cdot \sigma_y^2}            
                                                = \left( \dfrac{\sigma_{xy}}{\sigma_x \cdot \sigma_y} \right)^2   
                                                = r^2.
                              \end{aligned}\]
\end{frame}
\begin{frame}{Séries statistiques à deux variables}
\textbf{Variances et calcul pratique}

 Dans les formules ci-dessus, \(V(x)\) et \(V(y)\) désignent les variances de \(x\) et de \(y\) :
                        \[
                        V(x) = \sigma_x^2 = \dfrac{1}{n}\sum_{i=1}^{n}(x_i - \overline{x})^2 = \dfrac{1}{n}\sum_{i=1}^{n} x_i^2 - \overline{x}^2 ,
                        \qquad
                        V(y) = \sigma_y^2 = \dfrac{1}{n}\sum_{i=1}^{n} y_i^2 - \overline{y}^2 .
                        \]
\end{frame}
\begin{frame}{Séries statistiques à deux variables (2)}
En pratique, on calcule \(\sum x_i\), \(\sum y_i\), \(\sum x_i^2\), \(\sum y_i^2\) et \(\sum x_i y_i\) 
                        (dans un tableau ou à la calculatrice), puis les moyennes, les variances, la covariance,
                        les coefficients directeurs et \(r\). On admet que le coefficient \(r\) est toujours compris entre \(-1\) et \(1\).
\end{frame}
\begin{frame}{Séries statistiques à deux variables}
Pour la série des salaires et des candidats ci-dessus : \(\sum x_i^2 = 10\,590\) et 
                              \(\sum y_i^2 = 1\,360\), donc \(V(x) = \dfrac{10\,590}{5} - 46^2 = 2\) et \(V(y) = \dfrac{1\,360}{5} - 16^2 = 16\).
                              La droite de régression de \(y\) en \(x\) a pour coefficient directeur \(m = \dfrac{5,6}{2} = 2,8\) ; 
                              son équation est \(y - 16 = 2,8\,(x - 46)\), soit \(y = 2,8\,x - 112,8\).
                              Enfin \(r = \dfrac{5,6}{\sqrt{2}\times 4} \approx 0,99\) : la corrélation est forte et positive.
\end{frame}
\begin{frame}{Séries statistiques à deux variables}
\textbf{Ajustement après un changement de variable}

 Lorsque le nuage de points a une forme nettement courbe, un ajustement affine de \(y\) en \(x\) 
                        n'est pas adapté. On cherche alors une nouvelle variable pour laquelle le nuage est presque rectiligne,
                        on fait un ajustement affine sur cette nouvelle série, puis on revient aux variables de départ.
\end{frame}
\begin{frame}{Séries statistiques à deux variables (2)}
\begin{itemize}
\item \emph{Ajustement exponentiel.} Si \(y\) semble être multiplié par un même nombre 
                                    quand \(x\) augmente d'une même quantité (croissance ou décroissance << en pourcentage >>),
                                    on pose \(z = \ln y\) (les \(y_i\) étant strictement positifs). Si la droite de régression 
                                    de \(z\) en \(x\) est \(z = a x + b\), alors \(y = e^{ax + b} = B\, e^{ax}\) avec \(B = e^{b}\).

\end{itemize}
\end{frame}
\begin{frame}{Séries statistiques à deux variables (3)}
\begin{itemize}
\item \emph{Ajustement logarithmique.} Si \(y\) augmente d'une même quantité chaque fois que 
                                    \(x\) est multiplié par un même nombre, on pose \(t = \ln x\) (les \(x_i\) étant strictement positifs).
                                    Si la droite de régression de \(y\) en \(t\) est \(y = a t + b\), alors \(y = a \ln x + b\).

\end{itemize}
\end{frame}
\begin{frame}{Séries statistiques à deux variables (4)}
\begin{itemize}
\item D'autres changements de variable sont possibles, par exemple \(X = x^2\) ou \(Y = \dfrac{1}{y}\) ;
                                    l'énoncé indique en général celui qu'il faut utiliser.

\end{itemize}
\end{frame}
\begin{frame}{Séries statistiques à deux variables (5)}
On compare les coefficients de corrélation des deux séries pour juger quel ajustement est le meilleur.
\end{frame}
\begin{frame}{Séries statistiques à deux variables}
\textbf{Prévisions}

 Une fois l'ajustement choisi, son équation permet d'estimer une valeur de \(y\) pour une valeur de \(x\) donnée
                        (et inversement). On parle d'\emph{interpolation} si \(x\) est dans l'intervalle des valeurs observées
                        et d'\emph{extrapolation} sinon ; une extrapolation suppose que la tendance observée se poursuit,
                        ce qui n'est jamais certain.
\end{frame}
\begin{frame}{Séries statistiques à deux variables}
Une prévision n'a de sens que si l'ajustement est de bonne qualité, c'est-à-dire si \(|r|\) est
                        proche de \(1\) (un usage courant est d'exiger \(|r| \geq 0,87\), soit environ \(\dfrac{\sqrt{3}}{2}\)). 
                        Pour estimer \(y\) connaissant \(x\), on utilise la droite de régression de \(y\) en \(x\) ;
                        pour estimer \(x\) connaissant \(y\), on utilise de préférence celle de \(x\) en \(y\).
                        Les deux droites passent par le point moyen \(G\) et sont d'autant plus proches que \(|r|\) est proche de \(1\).
\end{frame}
\end{document}
